\tocsection{\textbf{11/4/26 Lecture 29:} Van Kampen's Theorem and Abstract Simplicial Complexes}

Here are some corollaries from the results proved last time.

\begin{lemma}[\textbf{Corollary 29.1}]
    If $K$ is a (finite) connected simplicial complex, then $\pi_1(\abs{K}, v)$ is finitely presented.
\end{lemma}

\begin{lemma}[\textbf{Corollary 29.2}]
    If $K(2)$ denotes the $2$-skeleton, that is, the union of the vertices, edges, and triangles of $K$, then
    \[ \pi_1(\abs{K}, v) \cong \pi_1(\abs{K(2)}, v) \]
\end{lemma}

\begin{lemma}[\textbf{Corollary 29.3}]
    If $J$ and $K$ are connected simplicial complexes that intersect at a point $v$, then
    \[ \pi_1(\abs{J \cup K}, v) \cong \pi_1(\abs{J}, v) * \pi_1(\abs{K}, v) \]
\end{lemma}
\begin{proof}{Proof:}
    Take maximal trees $T_1$ and $T_2$ in $J$ and $K$ respectively. Let $T_0 = T_1 \cup T_2$. This is a tree in $J \cup K$. Then, the generators are the edges in $J \setminus T_1$ ($\{g_{ij}\}$) and $K \setminus T_2$ ($\{h_{ij}\}$), and the relations come from $2$-simplices in $J$ ($\{r_p\}$) or $K$ ($\{s_q\}$). Then,
    \[ \pi_1(\abs{J \cup K}, v) = \brac{\{g_{ij}\} \cup \{h_{ij}\} \mid \{r_p\} \cup \{s_q\}} = \pi_1(\abs{J}) * \pi_1(\abs{K}) \]

    $\qed$
\end{proof}

\begin{example}{\textbf{Example:} Two Circles}
    \fig{0.7}{Lecture29Fig1}{}
    \[ \pi_1(J \cup K) = \Z * \Z = \pi_1(\abs{J}) * \pi_1(\abs{K}) \]
\end{example}

\begin{theorem}[\textbf{Theorem 29.4:} Van Kampen's Theorem]
    Suppose $J, K$ are simplicial complexes, $J \cap K$ is a simplicial complex, and all of these are path connected. Choose $v \in J \cap K$. Let $j : J \cap K \to J$ and $k : J \cap K \to K$ be the inclusion maps. Then, $\pi_1(\abs{J \cup K}, v)$ can be obtained from $\pi_1(\abs{J}, v) * \pi_1(\abs{K}, v)$ by adding the relations
    \[ j_{*}(z)k_{*}(z)^{-1} \]
    where $z$ runs over the generators of $\pi_1(\abs{J \cap K}, v)$.
\end{theorem}
\begin{proof}{Proof:}
    Choose a maximal tree $T_0$ of $J \cap K$. Extend to maximal trees $T_1, T_2$ in $J, K$ respectively. Note $T_1 \cup T_2$ is a maximal tree in $J \cup K$.\\

    Then, $\pi_1(\abs{J \cup K}, v)$ is generated by the edges in $J \cup K$ not in $T_1 \cup T_2$, which are
    \[ \{\tt{edges in } J \setminus T_1\} \cup \{\tt{edges in } K \setminus T_2\} \]
    but there are some redundancies for edges in $(J \setminus T_1) \cap (K \setminus T_2)$, so we add the relations
    \[ a_{ij} = b_{ij}, \quad a_{ij} \in G(J, T_1),\ b_{ij} \in G(K, T_2) \]
    for edges $v_i v_j$ in $(J \cap K) \setminus T_0$.

    $\qed$
\end{proof}

\textbf{Special Case:} If $J, K$ are simply connected and $J \cap K$ is nonempty and path connected, then $\pi_1(J \cup K, v)$ is a quotient of
\[ \pi_1(J, v) * \pi_1(K, v) = \{e\} \]
so $J \cup K$ is simply connected.

\begin{example}{\textbf{Example:} The Projective Plane}
    Consider $P^2 = S^2/\Z_2$. We know $\pi_1(P^2) = \Z_2$ since $\pi_1(S^2)$ is trivial and $S^2 \to P^2$ is a covering map. We will show this with Van Kampen.\\

    \textbf{Claim:} $P^2 = \{\tt{M\"obius strip}\} \cup_{S^1} D^2$
    \fig{0.7}{Lecture29Fig2}{}
    The reason is that if
    \fig{0.6}{Lecture29Fig3}{}
    then elements in $E_1$ are identified with elements of $E_2$, so these can be represented by $E_1 \cong D^2$. For elements of the equator, these are locally identified by the M\"obius strip:
    \fig{0.6}{Lecture29Fig4}{}
    Then, let $J = $ M\"obius strip, $K = D^2$, $J \cap K = S^1$, and $J \cup K = P^2$. Thus, by Van Kampen, and from the fact that the inclusion of a generator of $\pi_1(S_1)$ into the M\"obious strip is a loop that runs around twice, we see
    \begin{align*}
        \pi_1(P^2) &= \frac{\pi_1(\tt{M\"obius strip}) * \pi_1(D^2)}{\brac{j_{*}(z)k_{*}(z)^{-1}}}\\
        &= \frac{\brac{z} * \brac{e}}{\brac{z^2}}\\
        &= \Z/2\Z
    \end{align*}
\end{example}

\begin{definition}[\textbf{Definition 29.5:} Finite Abstract Simplicial Complexes]
    A \emph{finite abstract simplicial complex} is a pair $(V, S)$ where $V$ is a finite set and $S$ is a collection of nonempty subsets of $V$ such that
    \begin{enumerate}[label = (\roman*)]
        \item $\{v\} \in S$ for all $v \in V$
        \item $\P(X) \setminus \{\emptyset\} \subseteq S$ for all $X \in S$
    \end{enumerate}
\end{definition}

\begin{lemma}[\textbf{Theorem 29.6}]
    Any finite abstract simplicial complex can be realized as a simplicial complex in $\R^n$ for some $n$.
\end{lemma}
\begin{proof}{Proof:}
    Say $V$ has $k$ elements. Let $\Delta$ be the $(k - 1)$-dimensional simplicial complex corresponding to
    \[ \left\{x \in \R^k \col x_i \geq 0,\ \sum_{i = 1}^{k} x_i = 1\right\} \]
    Any collection of vertices in $\Delta$ spans a simplex. Thus, if $V = \{v_1, \ldots, v_k\}$, match these to the vertices $\{e_1, \ldots, e_k\}$ of $\Delta$, so that for each $X \in S$, we see that the vertices corresponding to the elements of $X$ form a simplex. Do this for all $X \in S$ and we get a simplicial complex.

    $\qed$
\end{proof}

\begin{example}{\textbf{Example:} Realizing an Abstract Complex}
    $V = \{a, b, c\}$, $S = \{\{a\}, \{b\}, \{c\}, \{a, b\}, \{a, c\}\}$. Then, putting this in $\R^3$, we see the corresponding simplicial complex is
    \[ V' = \{e_1, e_2, e_3\}, \quad S' = \{\{e_1\}, \{e_2\}, \{e_3\}, \{e_1, e_2\}, \{e_1, e_3\}\} \]
    \fig{0.5}{Lecture29Fig5}{}
\end{example}
